\chapter{Dental Infection}

\section*{Definition and Overview}
A dental infection is an infection of the teeth, gums, or surrounding tissues, most commonly arising from untreated tooth decay or advanced gum disease. The most common form is a dental abscess, a collection of pus caused by bacteria that have reached the pulp of a tooth or the gum tissues. Dental infections cause pain, swelling, and fever, and they can spread to the face, neck, or bloodstream, where they become life-threatening. Prompt dental and medical treatment is essential. This chapter covers the causes, signs, and treatment of dental infections.

\section*{Epidemiology}
Dental infections are common and are a leading cause of emergency dental visits. Untreated decay is the most common cause, and access to dental care strongly influences who develops serious infections. Dental abscesses occur at all ages but are most common in children and young adults, and in people with poor oral hygiene, diabetes, or weakened immunity. While serious complications are rare, deep neck infections and bloodstream infections from dental sources are potentially fatal, and they are more common in people without timely access to care.

\section*{Aetiology and Pathophysiology}
\begin{itemize}
    \item \textbf{Tooth decay:} bacteria penetrate through enamel and dentine into the pulp, where infection and pus develop
    \item \textbf{Gum disease:} deep periodontal pockets allow bacteria to infect the tissues around the tooth root
    \item \textbf{Trauma:} cracked or injured teeth can allow bacteria to enter the pulp
    \item \textbf{Common organisms:} mixed mouth bacteria, including streptococci and anaerobes, cause the infection
    \item \textbf{Abscess formation:} pus collects at the tooth root or in the gum, and pressure causes severe pain
    \item \textbf{Spread:} infection can track into the facial tissues, the floor of the mouth, the neck, or the bloodstream
    \item \textbf{Risk factors:} poor oral hygiene, untreated decay, diabetes, smoking, and immunosuppression
\end{itemize}

\section*{Clinical Features}
\begin{itemize}
    \item Severe, throbbing toothache, often worse when lying down
    \item Swelling of the gum, cheek, face, or jaw
    \item Pain on biting, and a tooth that feels raised or loose
    \item Fever, chills, and a general feeling of being unwell
    \item Bad breath and an unpleasant taste if the abscess drains
    \item Difficulty opening the mouth, swallowing, or breathing in advanced infection
    \item Swelling under the eye or in the neck requires urgent care
\end{itemize}

\section*{Differential Diagnosis}
\begin{itemize}
    \item Dental abscess versus pulpitis (inflamed pulp without pus) or sinus pain, which can mimic toothache
    \item Periapical abscess versus periodontal abscess, which differ in origin and treatment
    \item Cellulitis, a spreading infection of the facial tissues without a localised collection
    \item Osteomyelitis, infection of the jawbone
    \item Salivary gland infection or other facial swellings
    \item Erupting wisdom teeth (pericoronitis) causing pain and swelling
\end{itemize}

\section*{When to Seek Medical Care}
See a dentist urgently for toothache with swelling, fever, or any sign of spreading infection. Dental infections do not resolve on their own, and delaying treatment allows them to spread. People without a regular dentist should see a GP or attend an emergency department.

\textbf{Call 000 (or local emergency services) for:}
\begin{itemize}
    \item Swelling that spreads to the face, under the eye, or into the neck
    \item Difficulty breathing, swallowing, or opening the mouth
    \item High fever with swelling, confusion, or a rapid heart rate (possible sepsis)
    \item Swelling that makes the airway feel tight
\end{itemize}

\section*{Diagnosis and Investigations}
\begin{itemize}
    \item \textbf{Dental examination:} identification of the affected tooth and the source of the infection
    \item \textbf{X-rays:} dental X-rays show decay, abscesses, and bone loss; panoramic or CT imaging may be needed for spread
    \item \textbf{Blood tests:} elevated white cell count and inflammatory markers in significant infection
    \item \textbf{Vitality testing:} assessment of whether the tooth's nerve is alive
    \item \textbf{Imaging of the face and neck:} CT for deep or spreading infections
\end{itemize}

\section*{Management}
\begin{itemize}
    \item \textbf{Drainage:} releasing the pus is the definitive treatment; this may be via root canal treatment, incision of the gum, or extraction of the tooth
    \item \textbf{Root canal treatment:} removing the infected pulp, cleaning the canals, and sealing the tooth
    \item \textbf{Antibiotics:} prescribed when there is spreading infection, swelling, fever, or systemic illness; they are not a substitute for drainage
    \item \textbf{Pain relief:} paracetamol and anti-inflammatory medicines
    \item \textbf{Hospital care:} intravenous antibiotics and surgical drainage for deep or spreading infections
    \item \textbf{Treatment of gum infections:} scaling and root planing for periodontal abscesses
    \item \textbf{Follow-up:} completion of treatment and review to prevent recurrence
\end{itemize}

\section*{Complications}
\begin{itemize}
    \item Spread of infection to the facial spaces, the neck, or the floor of the mouth, which can obstruct the airway
    \item Ludwig's angina, a severe infection of the floor of the mouth that is life-threatening
    \item Sepsis from bacteria entering the bloodstream
    \item Cavernous sinus thrombosis from spread of facial infection
    \item Loss of the tooth despite treatment
    \item Damage to nearby teeth, nerves, or sinuses
\end{itemize}

\section*{Prognosis and Outcomes}
With prompt drainage and appropriate treatment, the vast majority of dental infections resolve completely, and the tooth can often be saved with root canal treatment. Serious complications are uncommon when care is timely but can be fatal if treatment is delayed. Pain usually improves rapidly once the infection is drained. Good ongoing oral care prevents new decay and further infections, and the long-term outlook is excellent with regular dental care.

\section*{Prevention and Self-Care}
\begin{itemize}
    \item Brush twice daily with fluoride toothpaste and floss daily
    \item Limit sugary foods and drinks
    \item Attend regular dental check-ups so decay is treated early
    \item Seek dental care promptly for toothache, rather than waiting
    \item Manage diabetes and other conditions that increase infection risk
    \item Do not smoke
    \item Complete any prescribed course of antibiotics
\end{itemize}

\section*{Sources and Further Reading}
The following reputable sources provide further information for healthcare students and patients seeking reliable, current advice about dental infections.

\begin{itemize}
    \item Australian Dental Association. \textit{Dental abscess: signs, treatment, and prevention.} https://www.ada.org.au/
    \item NHS. \textit{Dental abscess: symptoms and treatment.} https://www.nhs.uk/conditions/dental-abscess/
    \item Mayo Clinic. \textit{Tooth abscess: symptoms and treatment.} https://www.mayoclinic.org/diseases-conditions/tooth-abscess/
    \item MSD Manual. \textit{Dental abscess and infections: professional overview.} https://www.msdmanuals.com/
    \item Healthdirect Australia. \textit{Dental infections and abscesses.} https://www.healthdirect.gov.au/
    \item MedlinePlus. \textit{Tooth abscess.} https://medlineplus.gov/
\end{itemize}
